\section{Finite regions and their algebras}

\begin{definition}[Finite-region local algebra]
    \label{def:qca_local_algebra}
    \lean{SpinChain.Config,SpinChain.LocalAlgebra}
    \leanok
    Fix a nonnegative integer $d$.  For a finite region
    $\Lambda\subset\mathbb Z$, the configuration space and local observable
    algebra are
    \begin{align}
        \operatorname{Config}_d(\Lambda)
         & =\{0,\ldots,d-1\}^{\Lambda},                     \\
        \mathcal A_\Lambda
         & =M_{\operatorname{Config}_d(\Lambda)}(\mathbb C)
        \cong M_d^{\otimes\Lambda}.
    \end{align}
    The algebra $\mathcal A_\Lambda$ carries its operator norm.  This is the
    finite-region construction used in the QCA appendix of
    \cite{Cirac2017MPU}; no completion is taken here.
\end{definition}

\begin{definition}[Canonical inclusion of finite-region observables]
    \label{def:qca_local_inclusion}
    \lean{SpinChain.localInclusion}
    \leanok
    \uses{def:qca_local_algebra}
    If $\Lambda\subseteq\Gamma$ are finite regions, define
    \begin{align}
        \iota_{\Lambda,\Gamma} & \colon\mathcal A_\Lambda\longrightarrow
        \mathcal A_\Gamma,                                                      \\
        \iota_{\Lambda,\Gamma}(A)
                               & = A\otimes \mathbf 1_{\Gamma\setminus\Lambda}.
    \end{align}
    Here product configurations are identified with configurations on
    $\Gamma$.  This is a unital star-algebra homomorphism.
\end{definition}

\begin{lemma}[Injectivity of the finite-region inclusion]
    \label{lem:qca_local_inclusion_injective}
    \lean{SpinChain.localInclusion_injective}
    \leanok
    \uses{def:qca_local_inclusion}
    For $d>0$, the map $\iota_{\Lambda,\Gamma}$ is injective.
\end{lemma}

\begin{proof}\leanok
    Since $d>0$, choose a configuration $c$ on
    $\Gamma\setminus\Lambda$.  For configurations $x,y$ on $\Lambda$,
    \begin{align}
        \iota_{\Lambda,\Gamma}(A)((x,c),(y,c))
         & =A(x,y)\cdot\mathbf 1_{\Gamma\setminus\Lambda}(c,c)
        =A(x,y).
    \end{align}
    Hence $\iota_{\Lambda,\Gamma}(A)=\iota_{\Lambda,\Gamma}(B)$ implies
    $A=B$.
\end{proof}

\begin{lemma}[Operator-norm preservation]
    \label{lem:qca_local_inclusion_norm}
    \lean{SpinChain.localInclusion_norm}
    \leanok
    \uses{def:qca_local_inclusion}
    If $d>0$, then for every $A\in\mathcal A_\Lambda$,
    \begin{align}
        \lVert\iota_{\Lambda,\Gamma}(A)\rVert & =\lVert A\rVert.
    \end{align}
\end{lemma}

\begin{proof}\leanok
    \uses{lem:qca_local_inclusion_injective}
    An injective star-algebra homomorphism between complex $C^*$-algebras is
    isometric.
\end{proof}

\begin{lemma}[Isometry of the finite-region inclusion]
    \label{lem:qca_local_inclusion_isometry}
    \lean{SpinChain.localInclusion_isometry}
    \leanok
    \uses{def:qca_local_inclusion}
    For $d>0$, the map $\iota_{\Lambda,\Gamma}$ is an operator-norm isometry.
\end{lemma}

\begin{proof}\leanok
    \uses{lem:qca_local_inclusion_injective}
    Injectivity and the general theorem that an injective star-algebra
    homomorphism between complex $C^*$-algebras is an isometry give the norm identity.
\end{proof}

\begin{lemma}[Identity inclusion]
    \label{lem:qca_local_inclusion_refl}
    \lean{SpinChain.localInclusion_refl}
    \leanok
    \uses{def:qca_local_inclusion}
    For every finite region $\Lambda$,
    $\iota_{\Lambda,\Lambda}$ is the identity on $\mathcal A_\Lambda$.
\end{lemma}

\begin{proof}\leanok
    When the two regions agree, their complement is empty.  Hence
    \begin{align}
        \iota_{\Lambda,\Lambda}(A)
         & =A\otimes\mathbf 1_{\varnothing}
        =A\otimes 1
        =A.
    \end{align}
\end{proof}

\begin{lemma}[Composition of finite-region inclusions]
    \label{lem:qca_local_inclusion_trans}
    \lean{SpinChain.localInclusion_trans}
    \leanok
    \uses{def:qca_local_inclusion}
    If $\Lambda\subseteq\Gamma\subseteq\Delta$, then
    \begin{align}
        \iota_{\Gamma,\Delta}\circ\iota_{\Lambda,\Gamma}
         & =\iota_{\Lambda,\Delta}.
    \end{align}
\end{lemma}

\begin{proof}\leanok
    The complement decomposes as the disjoint union
    \begin{align}
        \Delta\setminus\Lambda
         & =(\Gamma\setminus\Lambda)\sqcup(\Delta\setminus\Gamma), \\
        \mathbf 1_{\Delta\setminus\Lambda}
         & =\mathbf 1_{\Gamma\setminus\Lambda}
        \otimes\mathbf 1_{\Delta\setminus\Gamma}.
    \end{align}
    Therefore
    \begin{align}
        \iota_{\Gamma,\Delta}(\iota_{\Lambda,\Gamma}(A))
         & =(A\otimes\mathbf 1_{\Gamma\setminus\Lambda})
        \otimes\mathbf 1_{\Delta\setminus\Gamma}         \\
         & =A\otimes\mathbf 1_{\Delta\setminus\Lambda}   \\
         & =\iota_{\Lambda,\Delta}(A).
    \end{align}
\end{proof}

\section{Translations of finite regions}

\begin{definition}[Translation of a finite region]
    \label{def:qca_translate_region}
    \lean{SpinChain.translateRegion,SpinChain.mem_translateRegion,
        SpinChain.mem_translateRegion_add,SpinChain.translateRegion_zero,
        SpinChain.translateRegion_add,
        SpinChain.translateRegion_mono}
    \leanok
    For a finite region $\Lambda\subset\mathbb Z$ and $a\in\mathbb Z$, define
    \begin{align}
        \Lambda+a & =\{i+a\mid i\in\Lambda\}.
    \end{align}
    Thus
    \begin{align}
        j\in\Lambda+a & \iff j-a\in\Lambda, \\
        \Lambda+0     & =\Lambda,           \\
        (\Lambda+a)+b & =\Lambda+(a+b).
    \end{align}
    If $\Lambda\subseteq\Gamma$, then $\Lambda+a\subseteq\Gamma+a$.
\end{definition}

\begin{definition}[Sumset of finite regions]
    \label{def:qca_region_sumset}
    \lean{SpinChain.regionSumset,SpinChain.mem_regionSumset,
        SpinChain.regionSumset_singleton_right,SpinChain.regionSumset_mono,
        SpinChain.regionSumset_mono_left,SpinChain.regionSumset_mono_right,
        SpinChain.regionSumset_zero_left,SpinChain.regionSumset_zero_right,
        SpinChain.regionSumset_assoc}
    \leanok
    \uses{def:qca_translate_region}
    For finite regions $\Lambda,\mathcal N\subset\mathbb Z$, define their
    sumset by
    \begin{align}
        \Lambda+\mathcal N
         & =\{i+a\mid i\in\Lambda,\ a\in\mathcal N\}.
    \end{align}
    Hence $j\in\Lambda+\mathcal N$ exactly when $j=i+a$ for some
    $i\in\Lambda$ and $a\in\mathcal N$. The sumset is monotone in both
    arguments, and the singleton $\{0\}$ is an identity. A singleton
    displacement recovers translation,
    \begin{align}
        \Lambda+\{a\} & =\Lambda+a.
    \end{align}
    Sumsets are associative:
    \begin{align}
        (\Lambda+\mathcal N)+\mathcal M & =\Lambda+(\mathcal N+\mathcal M).
    \end{align}
    These identities support composition of propagation bounds; in the QCA appendix of \cite{Cirac2017MPU}, $\Lambda+\mathcal N$ is the
    enlarged finite region in the propagation
    condition.
\end{definition}

\begin{definition}[Translation of sites and configurations]
    \label{def:qca_finite_configuration_translation}
    \lean{SpinChain.siteTranslation_apply_val,
        SpinChain.siteTranslation_symm_apply_val,SpinChain.Config.translation,
        SpinChain.Config.translation_apply,SpinChain.Config.translation_symm_apply}
    \leanok
    \uses{def:qca_local_algebra, def:qca_translate_region}
    Translation gives bijections
    \begin{align}
        t_{a,\Lambda} & \colon\Lambda\longrightarrow\Lambda+a,
                      & t_{a,\Lambda}(i)                       & =i+a,    \\
        T_{a,\Lambda} & \colon\operatorname{Config}_d(\Lambda)
        \longrightarrow\operatorname{Config}_d(\Lambda+a),
                      & (T_{a,\Lambda}x)(j)                    & =x(j-a).
    \end{align}
    Their inverses translate sites and configurations by $-a$.
\end{definition}

\begin{definition}[Finite-region translation of local observables]
    \label{def:qca_finite_local_translation}
    \lean{SpinChain.localTranslation,SpinChain.localTranslation_apply}
    \leanok
    \uses{def:qca_local_algebra, def:qca_finite_configuration_translation}
    Relabelling configurations defines a star-algebra equivalence
    \begin{align}
        \tau_{a,\Lambda} & \colon\mathcal A_\Lambda
        \longrightarrow\mathcal A_{\Lambda+a},                          \\
        \tau_{a,\Lambda}(A)(x,y)
                         & =A(T_{a,\Lambda}^{-1}x,T_{a,\Lambda}^{-1}y).
    \end{align}
    This is the finite-region translation associated with the lattice
    translation $i\mapsto i+a$.
\end{definition}

\begin{lemma}[Norm preservation of finite-region translation]
    \label{lem:qca_finite_local_translation_norm}
    \lean{SpinChain.localTranslation_norm,
        SpinChain.localTranslation_isometry}
    \leanok
    \uses{def:qca_finite_local_translation}
    Let $d$ be a nonnegative integer, let $\Lambda\subset\mathbb Z$ be finite,
    and let $a\in\mathbb Z$. For every $A\in\mathcal A_\Lambda$,
    \begin{align}
        \lVert\tau_{a,\Lambda}(A)\rVert & =\lVert A\rVert.
    \end{align}
    Consequently, $\tau_{a,\Lambda}$ is an isometry from
    $\mathcal A_\Lambda$ onto $\mathcal A_{\Lambda+a}$.
\end{lemma}

\begin{proof}\leanok
    \uses{def:qca_finite_local_translation}
    Translation bijectively relabels the configuration basis. The resulting
    star-algebra equivalence between the two finite matrix algebras preserves
    the operator norm. Applying this identity to $A-B$ and using linearity gives
    \begin{align}
        \lVert\tau_{a,\Lambda}(A)-\tau_{a,\Lambda}(B)\rVert
         & =\lVert\tau_{a,\Lambda}(A-B)\rVert
        =\lVert A-B\rVert,
    \end{align}
    so distances are preserved.
\end{proof}

\begin{lemma}[Finite-region translation action laws]
    \label{lem:qca_finite_translation_action}
    \lean{SpinChain.siteTranslation_zero,SpinChain.siteTranslation_add,
        SpinChain.Config.translation_apply_siteTranslation,
        SpinChain.Config.translation_zero,SpinChain.Config.translation_add,
        SpinChain.localTranslation_apply_translation,
        SpinChain.localTranslation_zero,SpinChain.localTranslation_add}
    \leanok
    \uses{def:qca_finite_configuration_translation, def:qca_finite_local_translation}
    Translation by zero acts trivially on sites, configurations, and
    finite-region observables. Successive translations by $a$ and $b$ agree
    with translation by $a+b$. Pointwise, for $i\in\Lambda$,
    $x,y\in\operatorname{Config}_d(\Lambda)$, and
    $A\in\mathcal A_\Lambda$,
    \begin{align}
        t_{0,\Lambda}(i)                   & =i,                                      \\
        t_{b,\Lambda+a}(t_{a,\Lambda}(i))
                                           & =t_{a+b,\Lambda}(i),                     \\
        (T_{a,\Lambda}x)(t_{a,\Lambda}(i)) & =x(i),                                   \\
        (T_{b,\Lambda+a}T_{a,\Lambda}x)
        (t_{b,\Lambda+a}(t_{a,\Lambda}(i)))
                                           & =(T_{a+b,\Lambda}x)(t_{a+b,\Lambda}(i)), \\
        (\tau_{a,\Lambda}(A))(T_{a,\Lambda}x,T_{a,\Lambda}y)
                                           & =A(x,y),                                 \\
        (\tau_{b,\Lambda+a}\tau_{a,\Lambda}(A))
        (T_{b,\Lambda+a}T_{a,\Lambda}x,
        T_{b,\Lambda+a}T_{a,\Lambda}y)
                                           & =(\tau_{a+b,\Lambda}(A))
        (T_{a+b,\Lambda}x,T_{a+b,\Lambda}y).
    \end{align}
    These identities express the additive action without identifying the
    distinct configuration spaces before applying their canonical translation
    bijections.
\end{lemma}

\begin{proof}\leanok
    The site identities follow from
    \begin{align}
        (i+a)+b & =i+(a+b).
    \end{align}
    Relabelling a configuration is precomposition by the inverse site
    translation, so evaluation at the corresponding translated site gives
    \begin{align}
        (T_{a,\Lambda}x)(t_{a,\Lambda}(i)) & =x(i).
    \end{align}
    Applying this identity successively proves the configuration action laws.
    Finally, the local observable translation relabels both matrix indices by
    $T_{a,\Lambda}^{-1}$; evaluating on translated configurations therefore
    gives $A(x,y)$, and applying this entrywise identity successively proves the
    local observable action laws.
\end{proof}

\begin{lemma}[Naturality of finite-region translation]
    \label{lem:qca_finite_local_translation_naturality}
    \lean{SpinChain.localTranslation_localInclusion}
    \leanok
    \uses{def:qca_local_inclusion, def:qca_translate_region, def:qca_finite_local_translation}
    If $\Lambda\subseteq\Gamma$, then for every
    $A\in\mathcal A_\Lambda$,
    \begin{align}
        \tau_{a,\Gamma}(\iota_{\Lambda,\Gamma}(A))
         & =\iota_{\Lambda+a,\Gamma+a}(\tau_{a,\Lambda}(A)).
    \end{align}
\end{lemma}

\begin{proof}\leanok
    The lattice translation identifies $\Gamma\setminus\Lambda$ with
    $(\Gamma+a)\setminus(\Lambda+a)$. Hence relabelling configurations
    carries the identity on the first complement to the identity on the
    second complement, and
    \begin{align}
        \tau_{a,\Gamma}(A\otimes\mathbf 1_{\Gamma\setminus\Lambda})
         & =\tau_{a,\Lambda}(A)
        \otimes\mathbf 1_{(\Gamma+a)\setminus(\Lambda+a)}.
    \end{align}
\end{proof}

\section{The local and quasi-local algebras}

\begin{definition}[Algebraic local algebra]
    \label{def:qca_algebraic_local_algebra}
    \lean{SpinChain.instDirectedSystemLocalAlgebra,
        SpinChain.AlgebraicLocalAlgebra,SpinChain.localObservable,
        SpinChain.localObservable_localInclusion}
    \leanok
    \uses{def:qca_local_algebra, def:qca_local_inclusion, lem:qca_local_inclusion_refl, lem:qca_local_inclusion_trans}
    The finite-region local algebras form a directed system under the maps
    $\iota_{\Lambda,\Gamma}$.  Its algebraic direct limit is the algebraic
    local algebra
    \begin{align}
        \mathcal A_{\mathrm{loc}}
         & =\varinjlim_{\Lambda\Subset\mathbb Z}\mathcal A_\Lambda.
    \end{align}
    For every finite region $\Lambda$, write
    $\iota_\Lambda\colon\mathcal A_\Lambda\to\mathcal A_{\mathrm{loc}}$
    for the canonical inclusion.  Whenever $\Lambda\subseteq\Gamma$, these
    inclusions satisfy
    \begin{align}
        \iota_\Gamma\circ\iota_{\Lambda,\Gamma} & =\iota_\Lambda.
    \end{align}
    This is the union
    $\bigcup_{\Lambda\Subset\mathbb Z}\mathcal A_\Lambda$ in the Appendix of
    \cite{Cirac2017MPU}.  No norm completion is taken.
\end{definition}

\begin{definition}[Operator norm on the local algebra]
    \label{def:qca_algebraic_local_operator_norm}
    \lean{SpinChain.algebraicLocalOperatorNorm,
        SpinChain.algebraicLocalOperatorNorm_localObservable,
        SpinChain.instNormAlgebraicLocalAlgebra,
        SpinChain.norm_localObservable}
    \leanok
    \uses{def:qca_algebraic_local_algebra, lem:qca_local_inclusion_norm}
    Assume $d>0$.  If a local observable $A\in\mathcal A_{\mathrm{loc}}$
    is represented by $A_\Lambda\in\mathcal A_\Lambda$ on a finite region
    $\Lambda$, define its operator norm by
    \begin{align}
        \lVert A\rVert & =\lVert A_\Lambda\rVert.
    \end{align}
    This definition is independent of the finite region and chosen
    representative, giving the norm completed in the Appendix of
    \cite{Cirac2017MPU} to obtain the quasi-local algebra.
\end{definition}

\begin{proof}\leanok
    If $A_\Lambda$ and $A_\Gamma$ represent the same local observable, then
    the defining relation for the inductive limit gives a finite region
    $\Delta\supseteq\Lambda\cup\Gamma$ such that
    \begin{align}
        \iota_{\Lambda,\Delta}(A_\Lambda)
         & =\iota_{\Gamma,\Delta}(A_\Gamma).
    \end{align}
    Both inclusions preserve the operator norm, so
    $\lVert A_\Lambda\rVert=\lVert A_\Gamma\rVert$.
\end{proof}

\begin{definition}[Normed star algebra of local observables]
    \label{def:qca_algebraic_local_normed_star_algebra}
    \lean{SpinChain.instNormedAddCommGroupAlgebraicLocalAlgebra,
        SpinChain.instNormedRingAlgebraicLocalAlgebra,
        SpinChain.instNormedAlgebraAlgebraicLocalAlgebra}
    \leanok
    \uses{def:qca_algebraic_local_operator_norm}
    For $d>0$, the local algebra with the operator norm is a normed complex
    star algebra.  No completeness assertion is made.
\end{definition}

\begin{proof}\leanok
    Taking finite-region representatives of the local observables and enlarging
    their regions to a common finite region reduces the claims to a single matrix algebra: the norm axioms, the product
    inequality, and compatibility with complex scalar multiplication then
    follow from the corresponding properties of the finite-region matrix
    algebra.
\end{proof}

\begin{lemma}[C-star identity for local observables]
    \label{lem:qca_algebraic_local_norm_star_mul_self}
    \lean{SpinChain.norm_star_mul_self}
    \leanok
    \uses{def:qca_algebraic_local_normed_star_algebra}
    For $d>0$ and every local observable $A$,
    \begin{align}
        \lVert A^*A\rVert & =\lVert A\rVert^2.
    \end{align}
\end{lemma}

\begin{proof}\leanok
    For a finite region $\Lambda$ and a representative
    $A_\Lambda\in\mathcal A_\Lambda$,
    \begin{align}
        \lVert A^*A\rVert
         & =\lVert A_\Lambda^*A_\Lambda\rVert
        =\lVert A_\Lambda\rVert^2
        =\lVert A\rVert^2.
    \end{align}
\end{proof}

\begin{definition}[C-star axiom structure on local observables]
    \label{def:qca_algebraic_local_cstar_ring}
    \lean{SpinChain.instCStarRingAlgebraicLocalAlgebra}
    \leanok
    \uses{lem:qca_algebraic_local_norm_star_mul_self}
    The normed star algebra of local observables satisfies the $C^*$-norm identity.
\end{definition}

\begin{proof}\leanok
    The C-star identity gives the required norm inequality in the reverse
    direction; submultiplicativity gives the other direction.
\end{proof}

\begin{lemma}[Local observables embed isometrically]
    \label{lem:qca_local_observable_isometry}
    \lean{SpinChain.localObservable_isometry}
    \leanok
    \uses{def:qca_algebraic_local_operator_norm}
    For $d>0$ and every finite region $\Lambda$, the canonical map
    \begin{align}
        \iota_\Lambda\colon\mathcal A_\Lambda
         & \longrightarrow\mathcal A_{\mathrm{loc}}
    \end{align}
    is an operator-norm isometry.
\end{lemma}

\begin{proof}\leanok
    \uses{def:qca_algebraic_local_operator_norm}
    The norm of $\iota_\Lambda(A)$ is defined to be the operator norm of its
    representative $A\in\mathcal A_\Lambda$.  Hence
    $\lVert\iota_\Lambda(A)\rVert=\lVert A\rVert$; applying this equality to
    differences proves preservation of distance.
\end{proof}

\begin{lemma}[Star-algebra universal property]
    \label{thm:qca_algebraic_local_algebra_universal}
    \lean{SpinChain.algebraicLocalAlgebraLift,
        SpinChain.algebraicLocalAlgebraLift_localObservable,
        SpinChain.algebraicLocalAlgebra_starAlgHom_ext,
        SpinChain.algebraicLocalAlgebraLift_unique}
    \leanok
    \uses{def:qca_algebraic_local_algebra}
    Let $B$ be a complex star algebra.  Suppose that for every finite region
    $\Lambda$ there is a star-algebra homomorphism
    $g_\Lambda\colon\mathcal A_\Lambda\to B$.  Whenever
    $\Lambda\subseteq\Gamma$, suppose that
    \begin{align}
        g_\Gamma\circ\iota_{\Lambda,\Gamma} & =g_\Lambda.
    \end{align}
    Then there is a unique star-algebra homomorphism
    $g\colon\mathcal A_{\mathrm{loc}}\to B$ satisfying
    $g\circ\iota_\Lambda=g_\Lambda$ for every finite region $\Lambda$.
\end{lemma}

\begin{proof}\leanok
    Define $g$ on a local observable represented by
    $A\in\mathcal A_\Lambda$ as $g_\Lambda(A)$.  Compatibility makes this
    independent of the representing region.  Since each $g_\Lambda$ preserves
    addition, multiplication, the unit, scalar multiplication, and the star
    operation, so does $g$.  Every element of the direct limit has a
    finite-region representative, which also proves uniqueness.
\end{proof}

\begin{definition}[Support of a local observable]
    \label{def:qca_supported_in}
    \lean{SpinChain.SupportedIn}
    \leanok
    \uses{def:qca_algebraic_local_algebra}
    A local observable $A\in\mathcal A_{\mathrm{loc}}$ is supported in the
    finite region $\Lambda$ if it lies in the image of
    $\iota_\Lambda\colon\mathcal A_\Lambda\to\mathcal A_{\mathrm{loc}}$.
\end{definition}

\begin{lemma}[Finite supports of local observables]
    \label{thm:qca_local_observable_finite_support}
    \lean{SpinChain.localObservable_injective,SpinChain.SupportedIn.mono,
        SpinChain.SupportedIn.union_left,SpinChain.SupportedIn.union_right,
        SpinChain.SupportedIn.add,SpinChain.SupportedIn.mul,
        SpinChain.SupportedIn.star,SpinChain.SupportedIn.zero,
        SpinChain.SupportedIn.one,SpinChain.exists_supportedIn}
    \leanok
    \uses{def:qca_supported_in, lem:qca_local_inclusion_injective}
    Every local observable is supported in some finite region.  If $A$ is
    supported in $\Lambda$ and $\Lambda\subseteq\Gamma$, then $A$ is supported
    in $\Gamma$.  If $A$ and $B$ are supported in $\Lambda$ and $\Gamma$,
    respectively, then $A+B$ and $AB$ are supported in $\Lambda\cup\Gamma$.
    The adjoint $A^*$ is supported in $\Lambda$, and both $0$ and $1$ are
    supported in every finite region.  For positive on-site dimension, each
    map $\iota_\Lambda$ is injective.
\end{lemma}

\begin{proof}\leanok
    \uses{def:qca_algebraic_local_algebra, lem:qca_local_inclusion_injective}
    Enlarge representatives by tensoring with the identity on the added sites.
    Two representatives can therefore be moved to the finite union of their
    regions, where sums and products are formed.  The canonical inclusions are
    unital star-algebra homomorphisms, so they preserve the adjoint, zero, and
    the unit.  For $d>0$, every transition map between finite-region algebras is
    injective; the canonical map from each member of such a directed system to
    its direct limit is therefore injective.  Finally, every element of an
    algebraic direct limit has a representative in one member of the directed
    system.
\end{proof}

\begin{definition}[Algebraic-local translation homomorphism]
    \label{def:qca_algebraic_local_translation_hom}
    \lean{SpinChain.algebraicLocalTranslationHom,
        SpinChain.algebraicLocalTranslationHom_localObservable}
    \leanok
    \uses{def:qca_finite_local_translation, def:qca_algebraic_local_algebra}
    For $a\in\mathbb Z$, finite-region translations induce a star-algebra
    homomorphism
    \begin{align}
        \tau_a\colon\mathcal A_{\mathrm{loc}} & \longrightarrow
        \mathcal A_{\mathrm{loc}}
    \end{align}
    determined on every finite-region representative by
    \begin{align}
        \tau_a(\iota_\Lambda(A))
         & =\iota_{\Lambda+a}(\tau_{a,\Lambda}(A)).
    \end{align}
\end{definition}

\begin{proof}\leanok
    \uses{def:qca_translate_region, def:qca_algebraic_local_algebra, lem:qca_finite_local_translation_naturality, thm:qca_algebraic_local_algebra_universal}
    Naturality gives, for $\Lambda\subseteq\Gamma$,
    \begin{align}
        \iota_{\Gamma+a}\!\left(\tau_{a,\Gamma}(\iota_{\Lambda,\Gamma}(A))\right)
         & =\iota_{\Gamma+a}\!\left(\iota_{\Lambda+a,\Gamma+a}(\tau_{a,\Lambda}(A))\right) \\
         & =\iota_{\Lambda+a}(\tau_{a,\Lambda}(A)).
    \end{align}
    Hence the finite-region maps form a compatible family and induce the stated
    homomorphism on the algebraic direct limit.
\end{proof}

\begin{lemma}[Action laws for algebraic-local translation homomorphisms]
    \label{lem:qca_algebraic_local_translation_hom_action}
    \lean{SpinChain.algebraicLocalTranslationHom_zero,
        SpinChain.algebraicLocalTranslationHom_add}
    \leanok
    \uses{def:qca_algebraic_local_translation_hom}
    The homomorphisms satisfy
    \begin{align}
        \tau_0            & =\operatorname{id}, \\
        \tau_b\circ\tau_a & =\tau_{a+b}.
    \end{align}
\end{lemma}

\begin{proof}\leanok
    \uses{def:qca_translate_region, def:qca_algebraic_local_algebra, lem:qca_finite_translation_action, thm:qca_algebraic_local_algebra_universal}
    On a representative $A\in\mathcal A_\Lambda$, the two sides of the
    composition law are represented on the equal finite regions
    $(\Lambda+a)+b$ and $\Lambda+(a+b)$. The finite-region action identity
    therefore gives
    \begin{align}
        \tau_b(\tau_a(\iota_\Lambda(A)))
         & =\iota_{\Lambda+(a+b)}
        (\tau_{a+b,\Lambda}(A))
        =\tau_{a+b}(\iota_\Lambda(A)).
    \end{align}
    The zero law is proved in the same way. Equality on all finite-region
    representatives determines equality on the algebraic direct limit.
\end{proof}

\begin{definition}[Translation of algebraic-local observables]
    \label{def:qca_algebraic_local_translation}
    \lean{SpinChain.algebraicLocalTranslation,
        SpinChain.algebraicLocalTranslation_apply,
        SpinChain.algebraicLocalTranslation_localObservable}
    \leanok
    \uses{def:qca_algebraic_local_translation_hom}
    For every $a\in\mathbb Z$, the homomorphism $\tau_a$ is a star-algebra
    automorphism of $\mathcal A_{\mathrm{loc}}$, and on finite-region
    representatives
    \begin{align}
        \tau_a(\iota_\Lambda(A))
         & =\iota_{\Lambda+a}(\tau_{a,\Lambda}(A)).
    \end{align}
\end{definition}

\begin{proof}\leanok
    \uses{def:qca_algebraic_local_translation_hom, lem:qca_algebraic_local_translation_hom_action}
    The homomorphism action laws give
    \begin{align}
        \tau_{-a}\circ\tau_a & =\tau_0=\operatorname{id}, \\
        \tau_a\circ\tau_{-a} & =\tau_0=\operatorname{id}.
    \end{align}
    Thus $\tau_a$ is bijective with inverse $\tau_{-a}$. The formula on
    finite-region representatives is the defining formula for the homomorphism.
\end{proof}

\begin{lemma}[Algebraic-local translation action laws]
    \label{lem:qca_algebraic_local_translation_action}
    \lean{SpinChain.algebraicLocalTranslation_zero,
        SpinChain.algebraicLocalTranslation_add,
        SpinChain.algebraicLocalTranslation_symm}
    \leanok
    \uses{def:qca_algebraic_local_translation}
    The automorphisms form an additive action of $\mathbb Z$:
    \begin{align}
        \tau_0            & =\operatorname{id}, \\
        \tau_b\circ\tau_a & =\tau_{a+b},        \\
        \tau_a^{-1}       & =\tau_{-a}.
    \end{align}
\end{lemma}

\begin{proof}\leanok
    \uses{def:qca_algebraic_local_translation, lem:qca_algebraic_local_translation_hom_action}
    The first two identities are the homomorphism action laws expressed as
    identities of star-algebra automorphisms. Taking $b=-a$ in the composition
    law gives
    \begin{align}
        \tau_{-a}\circ\tau_a & =\tau_0=\operatorname{id}, \\
        \tau_a\circ\tau_{-a} & =\tau_0=\operatorname{id},
    \end{align}
    which is the inverse identity.
\end{proof}

\begin{lemma}[Exact translation of algebraic-local supports]
    \label{lem:qca_algebraic_local_translation_support}
    \lean{SpinChain.algebraicLocalTranslation_supportedIn,
        SpinChain.algebraicLocalTranslation_supportedIn_iff}
    \leanok
    \uses{def:qca_supported_in, def:qca_algebraic_local_translation}
    A local observable $A$ is supported in a finite region $\Lambda$ if and only
    if $\tau_a(A)$ is supported in $\Lambda+a$.
\end{lemma}

\begin{proof}\leanok
    \uses{def:qca_translate_region, def:qca_algebraic_local_translation, lem:qca_algebraic_local_translation_hom_action}
    If $A=\iota_\Lambda(A_\Lambda)$, then
    \begin{align}
        \tau_a(A)
         & =\iota_{\Lambda+a}(\tau_{a,\Lambda}(A_\Lambda)),
    \end{align}
    which proves the forward implication. Applying the same argument to
    $\tau_a(A)$ with displacement $-a$ and using
    $\tau_{-a}\tau_a=\operatorname{id}$ and
    $(\Lambda+a)-a=\Lambda$, for the reverse implication.
\end{proof}

\begin{lemma}[Norm isometry of algebraic-local translation]
    \label{lem:qca_algebraic_local_translation_isometry}
    \lean{SpinChain.norm_algebraicLocalTranslation,
        SpinChain.algebraicLocalTranslation_isometry}
    \leanok
    \uses{def:qca_algebraic_local_operator_norm, def:qca_algebraic_local_translation}
    If $d>0$, then every translation preserves the operator norm and distance:
    \begin{align}
        \lVert\tau_a(A)\rVert           & =\lVert A\rVert,   \\
        \lVert\tau_a(A)-\tau_a(B)\rVert & =\lVert A-B\rVert.
    \end{align}
\end{lemma}

\begin{proof}\leanok
    \uses{def:qca_algebraic_local_operator_norm, def:qca_algebraic_local_translation, lem:qca_finite_local_translation_norm, thm:qca_local_observable_finite_support}
    For $A_\Lambda\in\mathcal A_\Lambda$ with
    $A=\iota_\Lambda(A_\Lambda)$, finite-region translation is a star-algebra
    equivalence of matrix algebras, hence preserves the operator norm. Therefore
    \begin{align}
        \lVert\tau_a(A)\rVert
         & =\lVert\tau_{a,\Lambda}(A_\Lambda)\rVert
        =\lVert A_\Lambda\rVert
        =\lVert A\rVert.
    \end{align}
    Applying this equality to $A-B$ gives distance preservation.
\end{proof}

\begin{definition}[C-star structure on a completion]
    \label{def:completion_cstar_algebra}
    \lean{UniformSpace.Completion.toComplAlgHom,
        UniformSpace.Completion.toComplStarAlgHom,
        UniformSpace.Completion.coe_toComplStarAlgHom,
        UniformSpace.Completion.instStar,UniformSpace.Completion.coe_star,
        UniformSpace.Completion.instContinuousStar,
        UniformSpace.Completion.instStarRing,
        UniformSpace.Completion.instCStarRing,
        UniformSpace.Completion.instStarModule,
        UniformSpace.Completion.instCStarAlgebra}
    \leanok
    The norm completion of a complex normed star algebra satisfying the C-star
    identity carries the continuously extended involution and a C-star algebra
    structure. The canonical map into the completion is a star-algebra
    homomorphism.
\end{definition}

\begin{proof}\leanok
    Extend the isometric involution continuously. The algebraic identities and
    the C-star identity hold on the dense canonical image and hence on the
    completion.
\end{proof}

\begin{definition}[Quasi-local observable algebra]
    \label{def:qca_quasi_local_algebra}
    \lean{SpinChain.QuasiLocalAlgebra,SpinChain.algebraicToQuasiLocal}
    \leanok
    \uses{def:qca_algebraic_local_cstar_ring, def:qca_algebraic_local_normed_star_algebra, def:completion_cstar_algebra}
    For $d>0$, the quasi-local observable algebra is the operator-norm
    completion
    \begin{align}
        \mathcal A
         & =\overline{\mathcal A_{\mathrm{loc}}}^{\lVert\cdot\rVert}.
    \end{align}
    The involution on $\mathcal A$ is the unique continuous extension of the
    involution on $\mathcal A_{\mathrm{loc}}$.  It makes $\mathcal A$ a complex
    C-star algebra, and the canonical star-algebra homomorphism
    \begin{align}
        j\colon\mathcal A_{\mathrm{loc}} & \longrightarrow\mathcal A
    \end{align}
    is the completion map.
\end{definition}

\begin{proof}\leanok
    The involution is isometric on $\mathcal A_{\mathrm{loc}}$, hence extends
    continuously to the completion.  The identities
    $(A+B)^*=A^*+B^*$, $(AB)^*=B^*A^*$, and $(A^*)^*=A$
    hold on the dense subalgebra $j(\mathcal A_{\mathrm{loc}})$ and therefore
    on $\mathcal A$.  The conjugate-linear identity
    $(cA)^*=\overline cA^*$ extends by the same argument.  Density also extends
    \begin{align}
        \lVert A^*A\rVert & =\lVert A\rVert^2,
    \end{align}
    so the completed algebra satisfies the C-star identity.
\end{proof}

\begin{definition}[Finite-region quasi-local observable]
    \label{def:qca_quasi_local_observable}
    \lean{SpinChain.quasiLocalObservable}
    \leanok
    \uses{def:qca_quasi_local_algebra, def:qca_algebraic_local_algebra}
    For every finite region $\Lambda$, define
    \begin{align}
        j_\Lambda\colon\mathcal A_\Lambda & \longrightarrow\mathcal A,
                                          & j_\Lambda                  & =j\circ\iota_\Lambda.
    \end{align}
\end{definition}

\begin{lemma}[Finite-region observables in the quasi-local algebra]
    \label{lem:qca_quasi_local_observable}
    \lean{SpinChain.norm_algebraicToQuasiLocal,
        SpinChain.norm_quasiLocalObservable,
        SpinChain.quasiLocalObservable_isometry,
        SpinChain.quasiLocalObservable_injective,
        SpinChain.quasiLocalObservable_localInclusion}
    \leanok
    \uses{def:qca_quasi_local_observable, def:qca_algebraic_local_operator_norm, def:qca_algebraic_local_algebra}
    For every finite region $\Lambda$, the map $j_\Lambda$ is an injective
    isometry.  If $\Lambda\subseteq\Gamma$, then
    \begin{align}
        j_\Gamma\circ\iota_{\Lambda,\Gamma} & =j_\Lambda.
    \end{align}
\end{lemma}

\begin{proof}\leanok
    \uses{def:qca_quasi_local_observable, def:qca_algebraic_local_operator_norm, def:qca_algebraic_local_algebra}
    The algebraic inclusion $\iota_\Lambda$ and the completion map $j$ both
    preserve the norm.  Their composite is therefore an isometry and hence
    injective.  Compatibility follows from
    $\iota_\Gamma\circ\iota_{\Lambda,\Gamma}=\iota_\Lambda$.
\end{proof}

\begin{lemma}[Density of the canonical completion map]
    \label{thm:completion_map_dense_range}
    \lean{UniformSpace.Completion.denseRange_toComplStarAlgHom}
    \leanok
    \uses{def:completion_cstar_algebra}
    The canonical star-algebra homomorphism from a normed star algebra into its
    completion has dense range.
\end{lemma}

\begin{proof}\leanok
    This is the defining density property of the completion.
\end{proof}

\begin{lemma}[Density of finite-region observables]
    \label{thm:qca_quasi_local_density}
    \lean{SpinChain.denseRange_algebraicToQuasiLocal,
        SpinChain.dense_iUnion_range_quasiLocalObservable}
    \leanok
    \uses{def:qca_quasi_local_algebra, def:qca_quasi_local_observable, thm:qca_local_observable_finite_support, thm:completion_map_dense_range}
    Algebraic local observables are dense in the quasi-local algebra.  Equivalently,
    \begin{align}
        \overline{j(\mathcal A_{\mathrm{loc}})}
         & =\overline{\bigcup_{\Lambda\Subset\mathbb Z}
            j_\Lambda(\mathcal A_\Lambda)}
        =\mathcal A.
    \end{align}
\end{lemma}

\begin{proof}\leanok
    \uses{thm:qca_local_observable_finite_support}
    Density of $j(\mathcal A_{\mathrm{loc}})$ is the defining density property
    of the completion.  Every algebraic local observable has a representative
    in some finite region, so
    \begin{align}
        j(\mathcal A_{\mathrm{loc}})
         & =\bigcup_{\Lambda\Subset\mathbb Z}j_\Lambda(\mathcal A_\Lambda),
    \end{align}
    and taking closures yields
    \begin{align}
        \overline{\bigcup_{\Lambda\Subset\mathbb Z} j_\Lambda(\mathcal A_\Lambda)}
         & = \overline{j(\mathcal A_{\mathrm{loc}})} = \mathcal A.
    \end{align}
\end{proof}

In the setting of \cite[lines~2292--2306]{Cirac2017MPU}, lattice translation acts covariantly on the quasi-local
algebra.

\begin{definition}[Extension of star-algebra equivalences to completions]
    \label{def:completion_star_alg_equiv}
    \lean{UniformSpace.Completion.mapStarAlgEquiv,
        UniformSpace.Completion.mapStarAlgEquiv_apply,
        UniformSpace.Completion.mapStarAlgEquiv_coe}
    \leanok
    Let $A$ and $B$ be complex normed star algebras, and let
    $f\colon A\to B$ be a star-algebra equivalence.  If $f$ and $f^{-1}$ are
    continuous, then $f$ extends to a star-algebra equivalence
    \begin{align}
        \widehat f\colon\widehat A & \longrightarrow\widehat B
    \end{align}
    between their norm completions.  Its underlying map is the continuous
    extension of $f$.  Writing $j_A$ and $j_B$ for the canonical maps into the
    completions, the extension satisfies
    \begin{align}
        \widehat f(j_A(x)) & =j_B(f(x)).
    \end{align}
\end{definition}

\begin{proof}\leanok
    Applying the completion functor to $f$ and $f^{-1}$ shows that on the dense canonical
    images their composites satisfy
    \begin{align}
        \widehat{f^{-1}}\bigl(\widehat f(j_A(x))\bigr) & =j_A(x), &
        \widehat f\bigl(\widehat{f^{-1}}(j_B(y))\bigr) & =j_B(y).
    \end{align}
    Continuity gives the two inverse identities on the completions.
    Multiplication, addition, scalar multiplication, and the star operation
    commute with the extension because the corresponding identities hold on
    the canonical images and both sides are continuous.
\end{proof}

\begin{definition}[Translation of quasi-local observables]
    \label{def:qca_quasi_local_translation}
    \lean{SpinChain.quasiLocalTranslation,
        SpinChain.quasiLocalTranslation_apply,
        SpinChain.quasiLocalTranslation_algebraicToQuasiLocal,
        SpinChain.quasiLocalTranslation_quasiLocalObservable}
    \leanok
    \uses{def:completion_star_alg_equiv, def:qca_algebraic_local_translation, def:qca_quasi_local_algebra, def:qca_quasi_local_observable, lem:qca_algebraic_local_translation_isometry}
    For $d>0$ and every $a\in\mathbb Z$, algebraic-local translation extends to
    a star-algebra automorphism
    \begin{align}
        \tau_a\colon\mathcal A & \longrightarrow\mathcal A
    \end{align}
    of the quasi-local algebra.  On the dense algebraic-local subalgebra and on
    finite-region observables it satisfies
    \begin{align}
        \tau_a(j(A)) & =j(\tau_a(A)),
        \label{eq:qca_trans_dense}                                  \\
        \tau_a(j_\Lambda(A_\Lambda))
                     & =j_{\Lambda+a}(\tau_{a,\Lambda}(A_\Lambda)).
        \label{eq:qca_trans_finite}
    \end{align}
\end{definition}

\begin{proof}\leanok
    \uses{def:completion_star_alg_equiv, def:qca_algebraic_local_translation, def:qca_quasi_local_algebra, def:qca_quasi_local_observable, lem:qca_algebraic_local_translation_action, lem:qca_algebraic_local_translation_isometry}
    Algebraic-local translation and its inverse translation by $-a$ are
    operator-norm isometries, hence continuous.  Extend this equivalence to the
    norm completion to obtain~(\ref{eq:qca_trans_dense}).
    Substituting $A=\iota_\Lambda(A_\Lambda)$ and applying the finite-region
    translation formula yields~(\ref{eq:qca_trans_finite}).
\end{proof}

\section{Grouping sites into blocks}

\begin{definition}[Expansion of a blocked region]
    \label{def:qca_expanded_region}
    \lean{SpinChain.expandedRegion,SpinChain.blockSiteEquiv}
    \leanok
    Let $L>0$ be a block length and let $\Lambda\subset\mathbb Z$ be finite.
    The expansion of $\Lambda$ is
    \begin{align}
        \widehat\Lambda_L
         & =\{Lx+r:x\in\Lambda,\ 0\leq r<L\}.
    \end{align}
    Every $z\in\widehat\Lambda_L$ has unique block coordinates $(x,r)$ with
    $x\in\Lambda$ and $0\leq r<L$.  Increasing $r$ gives the consecutive-site
    order within the block, matching in \cite[lines~2308 and 2313--2320]{Cirac2017MPU} the finite-region site
    grouping.
\end{definition}

\begin{lemma}[Coordinates and monotonicity of expanded regions]
    \label{thm:qca_expanded_region_coordinates}
    \lean{SpinChain.mem_expandedRegion,SpinChain.expandedRegion_mono,
        SpinChain.mem_expandedRegion_blockSite,
        SpinChain.blockSiteEquiv_apply_fst_val,
        SpinChain.blockSiteEquiv_apply_snd,
        SpinChain.blockSiteEquiv_symm_apply_val}
    \leanok
    \uses{def:qca_expanded_region}
    An integer $z$ belongs to $\widehat\Lambda_L$ exactly when its quotient by
    $L$ belongs to $\Lambda$.  The inverse block-coordinate map sends $(x,r)$
    to $Lx+r$, and $\Lambda\subseteq\Gamma$ implies
    $\widehat\Lambda_L\subseteq\widehat\Gamma_L$.
\end{lemma}

\begin{proof}\leanok
    \uses{def:qca_expanded_region}
    Under Euclidean division by the positive integer $L$, the quotient and
    residue are unique, the residue lies in $\{0,\ldots,L-1\}$, and taking the
    inverse coordinates gives $Lx+r$.  Monotonicity follows by retaining the
    same quotient and residue.
\end{proof}

\begin{definition}[Blocking equivalence for finite configurations]
    \label{def:qca_configuration_blocking}
    \lean{SpinChain.Config.blocking}
    \leanok
    \uses{def:qca_expanded_region}
    For on-site dimension $d$, configurations of $d^L$-level blocked spins on
    $\Lambda$ are equivalent to configurations of $d$-level spins on
    $\widehat\Lambda_L$.  The blocked label at $x$ is decoded into its ordered
    length-$L$ word, and its $r$th letter is placed at $Lx+r$.
\end{definition}

\begin{lemma}[Evaluation and restriction of blocked configurations]
    \label{thm:qca_configuration_blocking_formulas}
    \lean{SpinChain.Config.blocking_apply,
        SpinChain.Config.blocking_symm_apply,
        SpinChain.Config.restrict_blocking}
    \leanok
    \uses{def:qca_configuration_blocking, thm:qca_expanded_region_coordinates}
    Evaluation at $Lx+r$ reads the $r$th letter of the blocked spin at $x$.
    Conversely, encoding collects these letters in increasing residue order.
    At the configuration level, decoding commutes with restriction from
    $\Gamma$ to $\Lambda\subseteq\Gamma$.
\end{lemma}

\begin{proof}\leanok
    \uses{def:qca_configuration_blocking, thm:qca_expanded_region_coordinates}
    The two evaluation formulas are the inverse laws for the ordered word
    encoding.  Restriction preserves the quotient block coordinate and residue,
    so both sides have the same value at every expanded site.
\end{proof}

\begin{definition}[Finite local-algebra blocking equivalence]
    \label{def:qca_local_blocking}
    \lean{SpinChain.localBlocking}
    \leanok
    \uses{def:qca_local_algebra, def:qca_configuration_blocking}
    Reindexing matrix rows and columns by the configuration blocking
    equivalence gives a star-algebra equivalence
    \begin{align}
        \mathcal A^{(d^L)}_\Lambda
         & \simeq \mathcal A^{(d)}_{\widehat\Lambda_L}.
    \end{align}
    This is the finite-region equivalence underlying the algebraic-local
    blocking construction below.  No quasi-local completion, compatibility
    with translations, propagation statement, circuit representation, or
    conjugation of quantum cellular automata is asserted here.
\end{definition}

\begin{lemma}[Finite local-algebra blocking formulas]
    \label{thm:qca_local_blocking_formulas}
    \lean{SpinChain.localBlocking_apply,SpinChain.localBlocking_norm,
        SpinChain.localBlocking_isometry}
    \leanok
    \uses{def:qca_local_blocking}
    Matrix entries are reindexed by applying the inverse configuration
    equivalence to both indices.  The blocking equivalence preserves the
    operator norm and is an isometry.
\end{lemma}

\begin{proof}\leanok
    \uses{def:qca_local_blocking}
    The entry formula is matrix reindexing.  A star-algebra equivalence between
    finite matrix algebras preserves the C-star norm and hence distances.
\end{proof}

\begin{definition}[Finite block hull]
    \label{def:qca_block_hull}
    \lean{SpinChain.blockHull}
    \leanok
    For $L>0$ and a finite original region $\Delta\subset\mathbb Z$, define
    its block hull by
    \begin{align}
        \operatorname{Hull}_L(\Delta)
         & =\{\lfloor z/L\rfloor:z\in\Delta\}.
    \end{align}
    This is the finite set of length-$L$ blocks that meet $\Delta$.
\end{definition}

\begin{lemma}[Block-hull containment and monotonicity]
    \label{lem:qca_block_hull_properties}
    \lean{SpinChain.subset_expandedRegion_blockHull,SpinChain.blockHull_mono,
        SpinChain.blockHull_expandedRegion}
    \leanok
    \uses{def:qca_block_hull, def:qca_expanded_region}
    Let $L>0$.  Every finite original region is contained in the expansion of
    its block hull,
    \begin{align}
        \Delta\subseteq
        \widehat{\operatorname{Hull}_L(\Delta)}_L.
    \end{align}
    If $\Delta\subseteq\Theta$, then
    $\operatorname{Hull}_L(\Delta)\subseteq
        \operatorname{Hull}_L(\Theta)$.  Expansion and block hull are exact in the
    blocked direction:
    \begin{align}
        \operatorname{Hull}_L(\widehat\Lambda_L) & =\Lambda.
    \end{align}
\end{lemma}

\begin{proof}\leanok
    \uses{def:qca_block_hull, thm:qca_expanded_region_coordinates}
    The quotient coordinate of each $z\in\Delta$ belongs to the image defining
    the hull, which proves containment after expansion.  Monotonicity follows
    by retaining the same point $z$ in the larger region.  Finally, every
    expanded site has quotient in $\Lambda$, while for each $x\in\Lambda$ the
    site $Lx$ belongs to $\widehat\Lambda_L$ and has quotient $x$.  Hence
    $\operatorname{Hull}_L(\widehat\Lambda_L)=\Lambda$.
\end{proof}

\begin{lemma}[Blocking commutes with finite local inclusions]
    \label{lem:qca_local_blocking_inclusion}
    \lean{SpinChain.localBlocking_localInclusion}
    \leanok
    \uses{def:qca_local_blocking, thm:qca_configuration_blocking_formulas, def:qca_local_inclusion}
    Let $L>0$.  For finite blocked regions $\Lambda\subseteq\Gamma$, the
    square formed by finite blocking and the two canonical local inclusions
    commutes:
    \begin{align}
        \mathcal B_{L,\Gamma}\,\iota_{\Lambda,\Gamma}
         & =\iota_{\widehat\Lambda_L,\widehat\Gamma_L}\,
        \mathcal B_{L,\Lambda}.
    \end{align}
\end{lemma}

\begin{proof}\leanok
    \uses{thm:qca_local_blocking_formulas, thm:qca_configuration_blocking_formulas}
    For configurations $x,y$ on $\widehat\Gamma_L$, let $\widetilde x$ and
    $\widetilde y$ be their block encodings on $\Gamma$.  The two matrix
    entries in the commuting square are
    \begin{align}
        \bigl[\mathcal B_{L,\Gamma}
            (\iota_{\Lambda,\Gamma}(A))\bigr]_{x,y}
         & =A_{\widetilde x|_\Lambda,\widetilde y|_\Lambda}
        \mathbf 1_{\{\widetilde x|_{\Gamma\setminus\Lambda}
        =\widetilde y|_{\Gamma\setminus\Lambda}\}},         \\
        \bigl[\iota_{\widehat\Lambda_L,\widehat\Gamma_L}
            (\mathcal B_{L,\Lambda}(A))\bigr]_{x,y}
         & =A_{\widetilde x|_\Lambda,\widetilde y|_\Lambda}
        \mathbf 1_{\{x|_{\widehat\Gamma_L\setminus\widehat\Lambda_L}
                =y|_{\widehat\Gamma_L\setminus\widehat\Lambda_L}\}}.
    \end{align}
    Restriction commutes with block encoding, and decoding identifies
    $\Gamma\setminus\Lambda$ with the original sites in
    $\widehat\Gamma_L\setminus\widehat\Lambda_L$.  Hence the two indicator
    conditions are equivalent and the entries agree.
\end{proof}

\begin{definition}[Forward and reverse algebraic-local blocking maps]
    \label{def:qca_algebraic_blocking_maps}
    \lean{SpinChain.algebraicLocalBlockingHom,
        SpinChain.algebraicLocalUnblockingHom}
    \leanok
    \uses{def:qca_algebraic_local_algebra, def:qca_local_blocking, def:qca_block_hull}
    Let $L>0$.  The compatible finite blocking maps induce a star-algebra
    homomorphism
    \begin{align}
        \mathcal B_L:\mathcal A_{\mathrm{loc}}^{(d^L)}
        \longrightarrow\mathcal A_{\mathrm{loc}}^{(d)}.
    \end{align}
    In the reverse direction there is a star-algebra homomorphism
    \begin{align}
        \mathcal U_L:\mathcal A_{\mathrm{loc}}^{(d)}
        \longrightarrow\mathcal A_{\mathrm{loc}}^{(d^L)}.
    \end{align}
    A representative on an arbitrary original region $\Delta$ is included into
    $\widehat{\operatorname{Hull}_L(\Delta)}_L$ before inverting the finite
    blocking.
\end{definition}

\begin{proof}\leanok
    \uses{lem:qca_block_hull_properties, lem:qca_local_blocking_inclusion, thm:qca_algebraic_local_algebra_universal}
    For $\Lambda\subseteq\Gamma$, the finite commuting square gives
    \begin{align}
        \iota_{\widehat\Gamma_L}
        \mathcal B_{L,\Gamma}(\iota_{\Lambda,\Gamma}(A))
         & =\iota_{\widehat\Lambda_L}\mathcal B_{L,\Lambda}(A),
    \end{align}
    so the forward finite-region maps form a compatible family.  For the
    reverse family, block-hull monotonicity sends $\Delta\subseteq\Theta$ to
    $\operatorname{Hull}_L(\Delta)\subseteq
        \operatorname{Hull}_L(\Theta)$.  Transitivity of the canonical inclusions
    and the inverse commuting square then give
    \begin{align}
        \iota_{\operatorname{Hull}_L(\Theta)}
        \mathcal U_{L,\Theta}(\iota_{\Delta,\Theta}(A))
         & =\iota_{\operatorname{Hull}_L(\Delta)}
        \mathcal U_{L,\Delta}(A).
    \end{align}
    The universal property of the algebraic direct limit therefore induces
    both star-algebra homomorphisms.
\end{proof}

\begin{lemma}[Evaluation of the algebraic-local blocking maps]
    \label{lem:qca_algebraic_blocking_maps_evaluation}
    \lean{SpinChain.algebraicLocalBlockingHom_localObservable,
        SpinChain.algebraicLocalUnblockingHom_localObservable,
        SpinChain.algebraicLocalBlockingHom_comp_unblockingHom}
    \leanok
    \uses{def:qca_algebraic_blocking_maps, lem:qca_block_hull_properties}
    Let $L>0$.  The forward map applies finite blocking to a representative on
    $\Lambda$ and regards the result as supported on $\widehat\Lambda_L$.  The
    reverse map enlarges a representative on $\Delta$ to the expansion of its
    block hull, applies the inverse finite blocking equivalence, and regards
    the result as supported on $\operatorname{Hull}_L(\Delta)$.  The forward
    composite is the identity,
    \begin{align}
        \mathcal B_L\mathcal U_L
         & =\operatorname{id}_{\mathcal A_{\mathrm{loc}}^{(d)}}.
    \end{align}
    This holds for every on-site dimension $d$.
\end{lemma}

\begin{proof}\leanok
    \uses{lem:qca_local_blocking_inclusion, lem:qca_block_hull_properties}
    These are the evaluation formulas for the two compatible-family lifts.
    Forward compatibility is the commuting square.  Reverse compatibility
    follows by monotonicity of block hulls, composition of local inclusions,
    and the same square applied in the inverse direction.  Forward blocking
    after the reverse map reduces to the canonical inclusion of $\Delta$ into
    its expanded block hull.
\end{proof}

\begin{definition}[Algebraic-local site-blocking equivalence]
    \label{def:qca_algebraic_blocking}
    \lean{SpinChain.algebraicLocalBlocking}
    \leanok
    \uses{def:qca_algebraic_blocking_maps}
    For $d>0$ and $L>0$, forward blocking gives a star-algebra equivalence
    \begin{align}
        \mathcal A_{\mathrm{loc}}^{(d^L)}
         & \simeq \mathcal A_{\mathrm{loc}}^{(d)}.
    \end{align}
    In \cite[lines~2308 and 2313--2320]{Cirac2017MPU}, this equivalence is the algebraic-local change of site
    grouping.
\end{definition}

\begin{proof}\leanok
    \uses{lem:qca_algebraic_blocking_maps_evaluation, thm:qca_local_observable_finite_support, lem:qca_local_blocking_inclusion}
    The identity $\mathcal B_L\mathcal U_L=\operatorname{id}$ makes
    $\mathcal B_L$ surjective.  For injectivity, represent two blocked
    observables on finite regions $\Lambda$ and $\Gamma$, enlarge both to
    $\Omega=\Lambda\cup\Gamma$, and suppose their images under
    $\mathcal B_L$ agree.  Injectivity of the canonical map from the finite
    algebra on $\widehat\Omega_L$, followed by injectivity of
    $\mathcal B_{L,\Omega}$ and the finite commuting square, gives equality of
    the two representatives in the algebra on $\Omega$.  Thus
    $\mathcal B_L$ is bijective, with inverse $\mathcal U_L$, and defines the
    stated star-algebra equivalence.
\end{proof}

\begin{lemma}[Algebraic-local blocking formulas and norm]
    \label{lem:qca_algebraic_blocking_formulas}
    \lean{SpinChain.algebraicLocalBlocking_apply,
        SpinChain.algebraicLocalBlocking_localObservable,
        SpinChain.algebraicLocalBlocking_symm_apply,
        SpinChain.algebraicLocalBlocking_symm_localObservable,
        SpinChain.algebraicLocalBlocking_symm_localObservable_expandedRegion,
        SpinChain.algebraicLocalBlocking_norm,
        SpinChain.algebraicLocalBlocking_isometry,
        SpinChain.algebraicLocalBlocking_symm_isometry}
    \leanok
    \uses{def:qca_algebraic_blocking, lem:qca_algebraic_blocking_maps_evaluation}
    Let $d>0$ and $L>0$.  The equivalence agrees with the forward lift, and
    its inverse agrees with the block-hull unblocking map.  For
    $A\in\mathcal A^{(d^L)}_\Lambda$ and
    $C\in\mathcal A^{(d)}_\Delta$,
    \begin{align}
        \mathcal B_L\bigl(\iota_\Lambda(A)\bigr)
         & =\iota_{\widehat\Lambda_L}
        \bigl(\mathcal B_{L,\Lambda}(A)\bigr),    \\
        \mathcal B_L^{-1}\bigl(\iota_\Delta(C)\bigr)
         & =\iota_{\operatorname{Hull}_L(\Delta)}
        \left(\mathcal B^{-1}_{L,\operatorname{Hull}_L(\Delta)}
        \bigl(\iota_{\Delta,
                    \widehat{\operatorname{Hull}_L(\Delta)}_L}(C)\bigr)\right).
    \end{align}
    In particular, if $C\in\mathcal A^{(d)}_{\widehat\Lambda_L}$, then
    \begin{align}
        \mathcal B_L^{-1}\bigl(\iota_{\widehat\Lambda_L}(C)\bigr)
         & =\iota_\Lambda\bigl(\mathcal B_{L,\Lambda}^{-1}(C)\bigr).
    \end{align}
    The forward and inverse maps are operator-norm isometries.  In particular,
    each preserves the norm.
\end{lemma}

\begin{proof}\leanok
    \uses{lem:qca_algebraic_blocking_maps_evaluation, lem:qca_block_hull_properties, thm:qca_local_blocking_formulas, def:qca_algebraic_local_operator_norm}
    The first two formulas are the evaluation identities for the forward and
    reverse maps.  For an expanded region, the block-hull identity removes the
    enlargement and leaves the inverse finite blocking equivalence.  On a
    finite representative, forward norm preservation reduces to norm
    preservation by finite blocking and by the canonical local inclusion;
    distance preservation follows by applying the norm identity to
    differences.  The inverse isometry follows from the forward isometry and
    the right-inverse identity, rather than repeating the norm calculation.
\end{proof}

\begin{definition}[Quasi-local site-blocking equivalence]
    \label{def:qca_quasi_local_blocking}
    \lean{SpinChain.quasiLocalBlocking}
    \leanok
    \uses{def:qca_quasi_local_algebra, def:qca_algebraic_blocking, lem:qca_algebraic_blocking_formulas}
    For $d>0$ and $L>0$, site blocking extends uniquely by continuity to a
    star-algebra equivalence
    \begin{align}
        \overline{\mathcal B}_L:
        \mathcal A^{(d^L)} & \simeq\mathcal A^{(d)}.
    \end{align}
    In \cite[lines~2308 and 2313--2320]{Cirac2017MPU}, this extension completes the change of site
    grouping.
\end{definition}

\begin{proof}\leanok
    \uses{lem:qca_algebraic_blocking_formulas, def:completion_star_alg_equiv}
    The algebraic blocking equivalence and its inverse are isometries, hence
    continuous.  Extending both maps to the operator-norm completions gives
    mutually inverse star-algebra homomorphisms.
\end{proof}

\begin{lemma}[Quasi-local blocking evaluation and norm]
    \label{lem:qca_quasi_local_blocking_formulas}
    \lean{SpinChain.quasiLocalBlocking_apply,
        SpinChain.quasiLocalBlocking_algebraicToQuasiLocal,
        SpinChain.quasiLocalBlocking_symm_algebraicToQuasiLocal,
        SpinChain.quasiLocalBlocking_quasiLocalObservable,
        SpinChain.quasiLocalBlocking_symm_quasiLocalObservable,
        SpinChain.quasiLocalBlocking_symm_quasiLocalObservable_expandedRegion,
        SpinChain.norm_quasiLocalBlocking,
        SpinChain.quasiLocalBlocking_isometry}
    \leanok
    \uses{def:qca_quasi_local_blocking, lem:qca_algebraic_blocking_formulas, lem:qca_quasi_local_observable}
    Let $d>0$ and $L>0$.  Write $j^{(d)}$ for the canonical map from the
    algebraic-local algebra into its completion.  The underlying map of
    $\overline{\mathcal B}_L$ is the continuous extension of $\mathcal B_L$,
    and on the two dense algebraic-local subalgebras,
    \begin{align}
        \overline{\mathcal B}_L\bigl(j^{(d^L)}(A)\bigr)
         & =j^{(d)}\bigl(\mathcal B_L(A)\bigr),
        \label{eq:qca_blocking_dense_fwd}              \\
        \overline{\mathcal B}_L^{-1}\bigl(j^{(d)}(C)\bigr)
         & =j^{(d^L)}\bigl(\mathcal B_L^{-1}(C)\bigr).
        \label{eq:qca_blocking_dense_rev}
    \end{align}
    For $A\in\mathcal A^{(d^L)}_\Lambda$ and
    $C\in\mathcal A^{(d)}_\Delta$, these identities become
    \begin{align}
        \overline{\mathcal B}_L\bigl(j^{(d^L)}_\Lambda(A)\bigr)
         & =j^{(d)}_{\widehat\Lambda_L}
        \bigl(\mathcal B_{L,\Lambda}(A)\bigr),
        \label{eq:qca_blocking_fin_fwd}               \\
        \overline{\mathcal B}_L^{-1}\bigl(j^{(d)}_\Delta(C)\bigr)
         & =j^{(d^L)}_{\operatorname{Hull}_L(\Delta)}
        \left(\mathcal B^{-1}_{L,\operatorname{Hull}_L(\Delta)}
        \bigl(\iota_{\Delta,
                    \widehat{\operatorname{Hull}_L(\Delta)}_L}(C)\bigr)\right).
        \label{eq:qca_blocking_fin_rev}
    \end{align}
    If $C\in\mathcal A^{(d)}_{\widehat\Lambda_L}$, the inverse formula reduces
    exactly to
    \begin{align}
        \overline{\mathcal B}_L^{-1}
        \bigl(j^{(d)}_{\widehat\Lambda_L}(C)\bigr)
         & =j^{(d^L)}_\Lambda\bigl(\mathcal B_{L,\Lambda}^{-1}(C)\bigr).
        \label{eq:qca_blocking_fin_exact}
    \end{align}
    Finally, for all quasi-local observables $X$ and $Y$ on the blocked chain,
    \begin{align}
        \lVert\overline{\mathcal B}_L(X)\rVert & =\lVert X\rVert,
        \label{eq:qca_blocking_norm}                                \\
        \lVert\overline{\mathcal B}_L(X)-
        \overline{\mathcal B}_L(Y)\rVert       & =\lVert X-Y\rVert.
        \label{eq:qca_blocking_dist}
    \end{align}
\end{lemma}

\begin{proof}\leanok
    \uses{def:completion_star_alg_equiv, def:qca_quasi_local_observable, lem:qca_algebraic_blocking_formulas, lem:qca_block_hull_properties}
    Equalities~(\ref{eq:qca_blocking_dense_fwd}) and~(\ref{eq:qca_blocking_dense_rev}) are the defining continuous extensions of the algebraic equivalence and its inverse.
    Composing them with the finite-region canonical maps gives~(\ref{eq:qca_blocking_fin_fwd}) and~(\ref{eq:qca_blocking_fin_rev}).
    The exact formula~(\ref{eq:qca_blocking_fin_exact}) uses $\operatorname{Hull}_L(\widehat\Lambda_L)=\Lambda$.
    Because a star-algebra equivalence of C-star algebras preserves the norm, (\ref{eq:qca_blocking_norm}) holds, and applying it to $X-Y$ gives~(\ref{eq:qca_blocking_dist}).
\end{proof}

\begin{lemma}[Expansion of a translated blocked region]
    \label{lem:qca_expanded_region_translation}
    \lean{SpinChain.expandedRegion_translateRegion}
    \leanok
    \uses{def:qca_translate_region, def:qca_expanded_region}
    Let $L>0$, let $a\in\mathbb Z$, and let $\Lambda\subset\mathbb Z$ be a
    finite blocked region.  Then
    \begin{align}
        \widehat{\Lambda+a}_L
         & =\widehat\Lambda_L+aL.
    \end{align}
    In \cite[lines~2308 and 2313--2320]{Cirac2017MPU}, this identity is the translation rule for the site
    grouping.
\end{lemma}

\begin{proof}\leanok
    \uses{thm:qca_expanded_region_coordinates}
    Each original site has the unique form $xL+r$ with $0\leq r<L$.
    Translation of the blocked coordinate by $a$ sends this site to
    $(x+a)L+r=xL+r+aL$.  The argument is unchanged when $x$ or $a$ is
    negative.
\end{proof}

\begin{definition}[Carry-aware blocked neighborhood]
    \label{def:qca_blocked_neighborhood}
    \lean{SpinChain.blockedNeighborhood}
    \leanok
    \uses{def:qca_expanded_region, def:qca_block_hull, def:qca_region_sumset}
    Let $L>0$ and let $\mathcal N\subset\mathbb Z$ be a finite original-lattice
    neighborhood.  Write $\widehat{\{0\}}_L=\{0,\ldots,L-1\}$ for the sites in
    the reference block.  Define the induced blocked-lattice neighborhood by
    \begin{align}
        \mathcal C_L(\mathcal N)
         & =\operatorname{Hull}_L
        \bigl(\widehat{\{0\}}_L+\mathcal N\bigr).
    \end{align}
    The full reference block is included because the block reached after a
    displacement can depend on the starting residue; in \cite[lines~2308 and 2313--2320]{Cirac2017MPU}, this set captures the finite-region
    geometry implicit in the grouping of
    sites.
\end{definition}

\begin{lemma}[Residue description of the blocked neighborhood]
    \label{lem:qca_blocked_neighborhood_membership}
    \lean{SpinChain.mem_blockedNeighborhood,
        SpinChain.mem_blockedNeighborhood_iff_exists_eq}
    \leanok
    \uses{def:qca_blocked_neighborhood, thm:qca_expanded_region_coordinates}
    Let $L>0$, let $\mathcal N\subset\mathbb Z$ be finite, and let
    $b\in\mathbb Z$.  Then
    \begin{align}
        b\in\mathcal C_L(\mathcal N)
        \quad\Longleftrightarrow\quad
         & \exists r\in\{0,\ldots,L-1\},\ \exists n\in\mathcal N,
        \quad \left\lfloor\frac{r+n}{L}\right\rfloor=b,             \\
        \quad\Longleftrightarrow\quad
         & \exists r,s\in\{0,\ldots,L-1\},\ \exists n\in\mathcal N,
        \quad r+n=bL+s.
    \end{align}
    The quotient is the Euclidean quotient by the positive integer $L$.
    Consequently these formulas include negative displacements and both
    positive and negative block carries.  In general,
    $\operatorname{Hull}_L(\mathcal N)$ alone omits some carries caused by the
    starting residue, making explicit for \cite[lines~2313--2320]{Cirac2017MPU} the implicit boundary
    geometry.
\end{lemma}

\begin{proof}\leanok
    \uses{def:qca_blocked_neighborhood, def:qca_block_hull, thm:qca_expanded_region_coordinates}
    An element of $\widehat{\{0\}}_L+\mathcal N$ has the form $r+n$ with
    $0\leq r<L$ and $n\in\mathcal N$.  Taking its block coordinate gives the
    first equivalence.  Euclidean division writes $r+n=bL+s$ with the unique
    remainder $0\leq s<L$, which gives the second equivalence without any sign
    restriction on $n$.
\end{proof}

\begin{lemma}[Exact block hull after adding a neighborhood]
    \label{lem:qca_block_hull_region_sumset}
    \lean{SpinChain.blockHull_regionSumset_expandedRegion}
    \leanok
    \uses{def:qca_blocked_neighborhood, def:qca_region_sumset, def:qca_expanded_region, def:qca_block_hull}
    Let $L>0$.  For every finite blocked region $\Lambda$ and every finite
    original-lattice neighborhood $\mathcal N$,
    \begin{align}
        \operatorname{Hull}_L
        \bigl(\widehat\Lambda_L+\mathcal N\bigr)
         & =\Lambda+\mathcal C_L(\mathcal N).
    \end{align}
    In \cite[lines~2313--2320]{Cirac2017MPU}, this equality is the exact finite-region carry identity associated with the site
    grouping.
\end{lemma}

\begin{proof}\leanok
    \uses{lem:qca_blocked_neighborhood_membership, thm:qca_expanded_region_coordinates}
    Every expanded site has the form $xL+r$ with $x\in\Lambda$ and $0\leq r<L$.
    For $n\in\mathcal N$, Euclidean division gives
    \begin{align}
        \left\lfloor\frac{xL+r+n}{L}\right\rfloor
         & =x+\left\lfloor\frac{r+n}{L}\right\rfloor.
    \end{align}
    The second summand ranges over exactly $\mathcal C_L(\mathcal N)$ by
    Lemma~\ref{lem:qca_blocked_neighborhood_membership}.  This proves both
    inclusions, including for negative $x$ and $n$.
\end{proof}

\begin{lemma}[Carry-aware expansion containment]
    \label{lem:qca_region_sumset_expanded_region_subset}
    \lean{SpinChain.regionSumset_expandedRegion_subset}
    \leanok
    \uses{lem:qca_block_hull_region_sumset, lem:qca_block_hull_properties}
    Let $L>0$.  For every finite blocked region $\Lambda$ and every finite
    original-lattice neighborhood $\mathcal N$,
    \begin{align}
        \widehat\Lambda_L+\mathcal N
         & \subseteq
        \widehat{\Lambda+\mathcal C_L(\mathcal N)}_L.
    \end{align}
\end{lemma}

\begin{proof}\leanok
    \uses{lem:qca_block_hull_region_sumset, lem:qca_block_hull_properties}
    Every finite original region is contained in the expansion of its block
    hull; applying this containment to
    $\widehat\Lambda_L+\mathcal N$ and using
    Lemma~\ref{lem:qca_block_hull_region_sumset} identifies that hull.
\end{proof}

\begin{lemma}[Algebraic-local blocking intertwines scaled translations]
    \label{lem:qca_algebraic_blocking_translation}
    \lean{SpinChain.algebraicLocalBlocking_translation}
    \leanok
    \uses{lem:qca_expanded_region_translation, def:qca_algebraic_blocking, def:qca_algebraic_local_translation}
    Let $d,L>0$.  If $\mathcal B_L$ denotes the algebraic-local blocking
    equivalence, then for every $a\in\mathbb Z$,
    \begin{align}
        \mathcal B_L\,\tau^{(d^L)}_a
         & =\tau^{(d)}_{aL}\,\mathcal B_L.
    \end{align}
    In \cite[lines~2308 and 2313--2320]{Cirac2017MPU}, this intertwining identity is the translation rule induced by the blocking of
    sites.
\end{lemma}

\begin{proof}\leanok
    \uses{lem:qca_expanded_region_translation, lem:qca_algebraic_blocking_formulas, def:qca_algebraic_local_translation}
    On a representative supported in a finite blocked region $\Lambda$, both
    sides apply the same matrix reindexing.  Their supports agree because
    \begin{align}
        \widehat{\Lambda+a}_L
         & =\widehat\Lambda_L+aL.
    \end{align}
    Equality on all finite-region representatives proves the identity on the
    algebraic direct limit.
\end{proof}

\begin{lemma}[Quasi-local blocking intertwines scaled translations]
    \label{lem:qca_quasi_local_blocking_translation}
    \lean{SpinChain.quasiLocalBlocking_translation}
    \leanok
    \uses{lem:qca_algebraic_blocking_translation, def:qca_quasi_local_blocking, def:qca_quasi_local_translation}
    Let $d,L>0$.  If $\overline{\mathcal B}_L$ denotes the quasi-local blocking
    equivalence, then for every $a\in\mathbb Z$,
    \begin{align}
        \overline{\mathcal B}_L\,\overline\tau^{(d^L)}_a
         & =\overline\tau^{(d)}_{aL}\,\overline{\mathcal B}_L.
    \end{align}
    In \cite[lines~2308 and 2313--2320]{Cirac2017MPU}, this identity is the completed translation rule induced by the site
    grouping.
\end{lemma}

\begin{proof}\leanok
    \uses{lem:qca_algebraic_blocking_translation, lem:qca_quasi_local_blocking_formulas, def:qca_quasi_local_translation}
    The two sides are continuous maps on the quasi-local algebra.  On the dense
    algebraic-local subalgebra their values agree by
    Lemma~\ref{lem:qca_algebraic_blocking_translation}.  They therefore agree
    on the completion.
\end{proof}

\section{States of the quasi-local algebra}

\begin{definition}[Coordinates on a consecutive integer interval]
    \label{def:local_interval_coordinates}
    \lean{SpinChain.intervalCoordinates, SpinChain.intervalCoordinates_apply}
    \leanok
    For $a\in\mathbb Z$ and $N\geq0$, put
    $I(a,N)=\{a,a+1,\ldots,a+N-1\}$. The site numbering $i\mapsto a+i$
    identifies configurations on $I(a,N)$ with length-$N$ configurations.
    Reindexing both matrix coordinates gives a star-algebra equivalence
    $C_{a,N}:\mathcal A_{I(a,N)}\to M_{d^N}(\mathbb C)$.
    Explicitly, $C_{a,N}(X)_{\sigma\tau}=X_{\widetilde\sigma\widetilde\tau}$,
    where $\widetilde\sigma(a+i)=\sigma_i$.
    This includes $N=0$ and the one-dimensional algebra on the empty interval.
\end{definition}

\begin{lemma}[Local inclusions in interval coordinates]
    \label{thm:local_interval_coordinate_inclusion}
    \lean{SpinChain.finiteChainRegion_subset_expanded,
        SpinChain.intervalCoordinates_localInclusion_apply}
    \leanok
    Let $b,c\geq0$, $I=I(a,N)$ and $J=I(a-b,b+N+c)$. Then $I\subseteq J$.
    For configurations $\sigma,\tau$ on $b+N+c$ consecutive sites,
    \begin{align}
        C_{a-b,b+N+c}(\iota_{I,J}(X))_{\sigma\tau}
         & =C_{a,N}(X)_{\sigma|_M\,\tau|_M}
        \prod_{i\notin M}\delta_{\sigma_i,\tau_i},
    \end{align}
    where $M=\{b,b+1,\ldots,b+N-1\}$ and $\sigma|_M$ denotes the
    configuration $i\mapsto\sigma_{b+i}$ on $N$ sites.
    In \cite[Appendix]{Cirac2017MPU}, this formula is the identity-tensor inclusion
    expressed in increasing integer-site coordinates.
\end{lemma}
\begin{proof}\leanok
    \uses{def:local_interval_coordinates, def:qca_local_inclusion}
    Restriction to $I$ selects the middle $N$ coordinates. Equality of the
    two configurations on $J\setminus I$ is equality on all added sites.
    Substituting these identifications into the local inclusion formula gives the matrix entry.
\end{proof}

\begin{theorem}[Consecutive observables in the quasi-local algebra]
    \label{thm:quasi_local_interval_observable}
    \lean{SpinChain.quasiLocalIntervalObservable,
        SpinChain.quasiLocalIntervalObservable_apply,
        SpinChain.norm_quasiLocalIntervalObservable}
    \leanok
    For $d>0$ and $a\in\mathbb Z$, coordinate reindexing and the finite-region
    inclusion give a star-algebra inclusion
    $j_{a,k}:M_{d^k}(\mathbb C)\to\mathcal A$, including $k=0$.
    It preserves the matrix operator norm: $\|j_{a,k}(X)\|=\|X\|$.
    In \cite[Section~3]{Nachtergaele1996Spectral}, this embedding is the interval observable
    convention.
\end{theorem}
\begin{proof}\leanok
    \uses{def:local_interval_coordinates, def:qca_quasi_local_observable}
    The coordinate equivalence is isometric between the finite matrix
    algebras, and the finite-region inclusion preserves norm in the completion.
\end{proof}

\begin{definition}[Compatible finite-region states]
    \label{def:qca_compatible_local_state}
    \lean{SpinChain.CompatibleLocalState,
        SpinChain.CompatibleLocalState.functional,
        SpinChain.CompatibleLocalState.compatible,
        SpinChain.CompatibleLocalState.norm_le,
        SpinChain.CompatibleLocalState.map_one,
        SpinChain.CompatibleLocalState.nonneg_star_mul_self}
    \leanok
    \uses{def:qca_local_algebra}
    A compatible family of finite-region states consists of complex-linear
    functionals $f_\Lambda\colon\mathcal A_\Lambda\to\C$, indexed by finite
    regions $\Lambda\Subset\Z$, such that, for every inclusion
    $\Lambda\subseteq\Gamma$ and every $X\in\mathcal A_\Lambda$,
    \begin{align}
        f_\Gamma(\iota_{\Lambda,\Gamma}(X)) & =f_\Lambda(X),
                                            & |f_\Lambda(X)|         & \leq\lVert X\rVert, \\
        f_\Lambda(I)                        & =1,
                                            & f_\Lambda(X^\dagger X) & \geq0.
    \end{align}
\end{definition}

\begin{definition}[Algebraic extension of compatible states]
    \label{def:qca_compatible_state_algebraic_extension}
    \lean{SpinChain.CompatibleLocalState.algebraicFunctional,
        SpinChain.CompatibleLocalState.algebraicContinuousFunctional}
    \leanok
    \uses{def:qca_compatible_local_state, def:qca_algebraic_local_algebra,
        def:qca_algebraic_local_normed_star_algebra}
    Compatibility determines a complex-linear functional
    $f_{\mathrm{loc}}\colon\mathcal A_{\mathrm{loc}}\to\C$ by
    $f_{\mathrm{loc}}(\iota_\Lambda(X))=f_\Lambda(X)$.
    For $d>0$, this functional is continuous in the operator-norm topology.
\end{definition}

\begin{lemma}[Contractivity of the algebraic extension]
    \label{thm:qca_compatible_state_algebraic_bound}
    \lean{SpinChain.CompatibleLocalState.algebraicFunctional_localObservable,
        SpinChain.CompatibleLocalState.algebraicFunctional_norm_le}
    \leanok
    \uses{def:qca_compatible_state_algebraic_extension}
    For every finite region $\Lambda$ and $X\in\mathcal A_\Lambda$,
    $f_{\mathrm{loc}}(\iota_\Lambda(X))=f_\Lambda(X)$.
    If $d>0$, then $|f_{\mathrm{loc}}(A)|\leq\lVert A\rVert$ for every
    $A\in\mathcal A_{\mathrm{loc}}$.
\end{lemma}

\begin{proof}
    \leanok
    \uses{def:qca_compatible_state_algebraic_extension,
        def:qca_algebraic_local_operator_norm}
    The direct-limit property gives the stated values. Every algebraic local
    observable has a finite-region representative, and its norm is the norm
    of that representative. Thus
    $|f_{\mathrm{loc}}(\iota_\Lambda(X))|
        =|f_\Lambda(X)|\leq\lVert X\rVert=\lVert\iota_\Lambda(X)\rVert$.
\end{proof}

\begin{definition}[Continuous extension to the quasi-local algebra]
    \label{def:qca_compatible_state_completion}
    \lean{SpinChain.CompatibleLocalState.quasiLocalFunctional}
    \leanok
    \uses{def:qca_quasi_local_algebra,
        thm:qca_compatible_state_algebraic_bound}
    For $d>0$, let $\omega\colon\mathcal A\to\C$ be the continuous linear
    extension of $f_{\mathrm{loc}}$ to the operator-norm completion.
\end{definition}

\begin{lemma}[Evaluation of the continuous extension]
    \label{thm:qca_compatible_state_evaluation}
    \lean{SpinChain.CompatibleLocalState.quasiLocalFunctional_algebraicToQuasiLocal,
        SpinChain.CompatibleLocalState.quasiLocalFunctional_quasiLocalObservable}
    \leanok
    \uses{def:qca_compatible_state_completion, def:qca_quasi_local_observable}
    For every algebraic local observable $A$, one has
    $\omega(j(A))=f_{\mathrm{loc}}(A)$. For every finite region $\Lambda$ and
    $X\in\mathcal A_\Lambda$, one has $\omega(j_\Lambda(X))=f_\Lambda(X)$.
\end{lemma}

\begin{proof}
    \leanok
    \uses{def:qca_compatible_state_completion,
        thm:qca_compatible_state_algebraic_bound}
    Continuous extension preserves the values on the original normed space,
    giving the finite-region formula at $A=\iota_\Lambda(X)$.
\end{proof}

\begin{lemma}[Unique positive normalized extension]
    \label{thm:qca_compatible_state_extension}
    \lean{SpinChain.CompatibleLocalState.quasiLocalFunctional_nonneg_star_mul_self,
        SpinChain.CompatibleLocalState.quasiLocalFunctional_one,
        SpinChain.CompatibleLocalState.quasiLocalFunctional_norm_le,
        SpinChain.CompatibleLocalState.norm_quasiLocalFunctional,
        SpinChain.CompatibleLocalState.quasiLocalFunctional_unique,
        SpinChain.CompatibleLocalState.existsUnique_quasiLocalFunctional}
    \leanok
    \uses{def:qca_compatible_state_completion}
    Every compatible family of finite-region states, with $d>0$, has a unique
    continuous complex-linear extension $\omega\colon\mathcal A\to\C$
    with the prescribed finite-region values. It satisfies
    $\omega(I)=1$, $\lVert\omega\rVert=1$, and, for every $A\in\mathcal A$,
    $|\omega(A)|\leq\lVert A\rVert$ and $\omega(A^\dagger A)\geq0$.
    Compatibility of the finite-region functionals is an explicit hypothesis.
\end{lemma}

\begin{proof}
    \leanok
    \uses{thm:qca_compatible_state_evaluation,
        thm:qca_compatible_state_algebraic_bound,
        def:qca_quasi_local_algebra,
        def:qca_algebraic_local_algebra}
    The sets defined by $|\omega(A)|\leq\lVert A\rVert$ and
    $\omega(A^\dagger A)\geq0$ are closed, since multiplication, involution,
    and $\omega$ are continuous. These inequalities hold on the dense
    algebraic local subalgebra, hence on its completion. Evaluation at a
    finite-region identity gives $\omega(I)=1$ and
    $1\leq\lVert\omega\rVert$; contractivity gives the reverse inequality.
    Two continuous extensions agree on all finite-region observables,
    then on the algebraic local algebra, and therefore on its completion.
\end{proof}

\begin{definition}[States and pure states of the quasi-local algebra]
    \label{def:qca_state_space_purity}
    \lean{SpinChain.quasiLocalStateSpace, SpinChain.IsPureQuasiLocalState}
    \leanok
    Let $d>0$, and let $\mathcal A_d$ denote the quasi-local observable algebra.
    Its state space consists of all continuous complex-linear functionals
    $\omega$ such that $\|\omega\|=1$, $\omega(I)=1$, and
    $\omega(X^*X)\ge0$ for every $X\in\mathcal A_d$.
    A state is pure if it is an extreme point of this set, viewed as a real
    convex set. Translation invariance is not required of these functionals.
    In \cite[Section~1 and Section~3]{Nachtergaele1996Spectral}, these conditions define states and
    purity.
\end{definition}

\begin{lemma}[Support in convex constituents]
    \label{thm:qca_state_convex_support}
    \lean{SpinChain.quasiLocalState_left_apply_eq_zero_of_convex_decomposition,
        SpinChain.quasiLocalState_left_projection_eq_one_of_convex_decomposition}
    \leanok
    Let $\phi,\psi$ be states, let $0<t<1$, and set
    $\omega=(1-t)\phi+t\psi$.
    If $\omega(X^*X)=0$, then $\phi(X^*X)=0$.
    If $\omega$ is a state, $P=P^*=P^2$, and $\omega(P)=1$, then
    $\phi(P)=1$.
\end{lemma}
\begin{proof}
    \leanok
    \uses{def:qca_state_space_purity}
    Positivity gives two nonnegative real numbers whose strictly positive
    weighted sum is zero, so both numbers vanish; applied to $I-P$,
    $(I-P)^*(I-P)=I-P$ and normalization gives the projection assertion.
\end{proof}

\begin{lemma}[Purity from unique supported-state determination]
    \label{thm:qca_purity_of_supported_uniqueness}
    \lean{SpinChain.isPureQuasiLocalState_of_unique_supported_state}
    \leanok
    Let $\omega$ be a state, and let $(P_i)_{i\in J}$ be any family of
    self-adjoint projections with $\omega(P_i)=1$.
    Suppose every state $\phi$ satisfying $\phi(P_i)=1$ for every $i\in J$
    equals $\omega$. Then $\omega$ is pure.
\end{lemma}
\begin{proof}
    \leanok
    \uses{thm:qca_state_convex_support, def:qca_state_space_purity}
    In a nontrivial convex decomposition of $\omega$, each constituent
    has expectation one on every $P_i$. The uniqueness hypothesis forces
    the left constituent to equal $\omega$, which is the extreme-point
    condition.
\end{proof}

\begin{lemma}[Convexity of the quasi-local state space]
    \label{thm:quasi_local_state_space_convex}
    \lean{SpinChain.convex_quasiLocalStateSpace}
    \leanok
    The normalized positive continuous functionals on the quasi-local algebra
    form a convex set over the real numbers.
\end{lemma}
\begin{proof}\leanok
    \uses{def:qca_state_space_purity}
    If $\varphi,\psi$ are states and $a,b\geq0$ satisfy $a+b=1$, then
    $\|a\varphi+b\psi\|\leq a+b=1$ and $(a\varphi+b\psi)(I)=1$.
    Since $\|I\|=1$, evaluation at the identity gives the reverse norm bound.
    Positivity follows by taking the same convex combination of their
    nonnegative values on adjoint squares.
\end{proof}

\begin{lemma}[Pure states in finite convex decompositions]
    \label{thm:finite_pure_state_decomposition}
    \lean{SpinChain.exists_eq_of_isPureQuasiLocalState_of_finite_decomposition}
    \leanok
    Let $\varphi$ be pure among all states of the quasi-local algebra.
    Suppose that $\omega_i$ are states indexed by a finite set, that
    $w_i\geq0$ and $\sum_iw_i=1$, and that
    $\varphi=\sum_iw_i\omega_i$. Then $\varphi=\omega_i$ for some $i$.
    No translation-invariance condition is imposed on these states.
\end{lemma}
\begin{proof}\leanok
    \uses{thm:quasi_local_state_space_convex, def:qca_state_space_purity}
    The state $\varphi$ belongs to the convex hull of the $\omega_i$.
    This hull is contained in the full state space, so extremality in the
    state space implies extremality in the smaller hull. Every extreme
    point of the convex hull of a finite set belongs to that set.
\end{proof}

\begin{lemma}[Expansion of a consecutive blocked interval]
    \label{lem:qca_blocking_interval_region}
    \lean{SpinChain.expandedRegion_finiteChainRegion}
    \leanok
    For $L>0$, $a\in\mathbb Z$ and $N\geq0$, the expansion of the
    blocked interval $I(a,N)$ is
    \begin{align}
        \widehat{I(a,N)}_L=I(aL,NL).
    \end{align}
    Negative starts and the empty interval are included.
\end{lemma}
\begin{proof}\leanok
    \uses{def:local_interval_coordinates, def:qca_expanded_region}
    Euclidean division by $L$ shows that the quotient $q=\lfloor z/L\rfloor$ of an integer $z$
    satisfies $a\leq q<a+N$ exactly when $aL\leq z<(a+N)L$.
\end{proof}

\begin{lemma}[Finite blocking in consecutive matrix coordinates]
    \label{thm:qca_blocking_interval_coordinates}
    \lean{SpinChain.intervalCoordinates_localBlocking}
    \leanok
    Let $L>0$, $d\geq0$, $a\in\mathbb Z$ and $N\geq0$.
    Let $e_{N,L}$ decode the $N$ blocked spins into $NL$ consecutive
    spins, placing residue $r$ of blocked site $i$ at site $Li+r$.
    Define the matrix reindexing
    $\mathcal R_{N,L}(X)_{\sigma\tau}
        =X_{e_{N,L}^{-1}(\sigma),e_{N,L}^{-1}(\tau)}$.
    Then
    \begin{align}
        C^{(d)}_{aL,NL}\,\mathcal B_L
        \bigl((C^{(d^L)}_{a,N})^{-1}(X)\bigr)
        =\mathcal R_{N,L}(X).
    \end{align}
    Here the target region is identified using
    $\widehat{I(a,N)}_L=I(aL,NL)$, expressing the site grouping of \cite[Appendix, lines~2308 and 2313--2320]{Cirac2017MPU}
    in consecutive interval
    coordinates.
\end{lemma}
\begin{proof}\leanok
    \uses{lem:qca_blocking_interval_region, def:local_interval_coordinates,
        def:qca_configuration_blocking, thm:qca_local_blocking_formulas}
    The site $L(a+i)+r$ has interval coordinate $Li+r$ relative to $aL$.
    Thus decoding followed by consecutive numbering agrees with flattening
    the blocked configurations on both matrix indices.
\end{proof}

\begin{lemma}[Blocking of quasi-local interval observables]
    \label{thm:qca_blocking_quasi_local_interval}
    \lean{SpinChain.quasiLocalBlocking_quasiLocalIntervalObservable}
    \leanok
    For $d>0$, $L>0$, $a\in\mathbb Z$, $N\geq0$ and
    $X\in M_{(d^L)^N}(\mathbb C)$,
    \begin{align}
        \overline{\mathcal B}_L
        \bigl(j^{(d^L)}_{a,N}(X)\bigr)
        =j^{(d)}_{aL,NL}\bigl(\mathcal R_{N,L}(X)\bigr).
    \end{align}
\end{lemma}
\begin{proof}\leanok
    \uses{thm:qca_blocking_interval_coordinates,
        lem:qca_quasi_local_blocking_formulas,
        thm:quasi_local_interval_observable}
    Evaluating the completed blocking equivalence on the finite interval
    representative and applying the finite coordinate identity gives the formula.
\end{proof}

\begin{definition}[Pushforward of functionals by a star-algebra equivalence]
    \label{def:qca_functional_star_equivalence}
    \lean{SpinChain.quasiLocalFunctionalCongr,
        SpinChain.quasiLocalFunctionalCongr_apply}
    \leanok
    Let $d_1,d_2>0$ and let
    $e:\mathcal A_{d_1}\to\mathcal A_{d_2}$ be a star-algebra isomorphism.
    The induced continuous complex-linear equivalence between functional
    spaces is $T_e\omega=\omega\circ e^{-1}$; thus
    $(T_e\omega)(X)=\omega(e^{-1}X)$.
\end{definition}

\begin{lemma}[Star-algebra equivalence preserves states and purity]
    \label{thm:qca_state_purity_star_equivalence}
    \lean{SpinChain.norm_quasiLocalFunctionalCongr,
        SpinChain.quasiLocalFunctionalCongr_mem_stateSpace_iff,
        SpinChain.quasiLocalFunctionalCongr_image_stateSpace,
        SpinChain.isPureQuasiLocalState_quasiLocalFunctionalCongr_iff}
    \leanok
    Under the preceding hypotheses, $\|T_e\omega\|=\|\omega\|$.
    In addition, $T_e$ maps the entire state space of $\mathcal A_{d_1}$ onto
    the entire state space of $\mathcal A_{d_2}$, and $T_e\omega$ is pure
    if and only if $\omega$ is pure. No translation invariance is assumed.
    Here purity follows the extreme-point convention of
    \cite[Section~3, lines~1469--1482]{Nachtergaele1996Spectral} on the state space.
\end{lemma}
\begin{proof}\leanok
    \uses{def:qca_functional_star_equivalence,def:qca_state_space_purity}
    A star-algebra isomorphism of complex C-star algebras is isometric.
    Inverse precomposition therefore preserves the dual norm. It also
    preserves the identity and adjoint squares, so it maps the state spaces
    onto one another. Its restriction to real scalars is a linear
    equivalence, and linear equivalences carry extreme points to extreme
    points.
\end{proof}

\begin{definition}[A blocked-chain functional on the original lattice]
    \label{def:qca_blocking_functional}
    \lean{SpinChain.quasiLocalBlockingFunctional,
        SpinChain.quasiLocalBlockingFunctional_apply}
    \leanok
    Let $d,L>0$ and let
    $\mathcal B_L:\mathcal A_{d^L}\to\mathcal A_d$ be the site-blocking
    isomorphism. A continuous functional $\omega$ on the blocked chain
    defines the original-lattice functional
    $T_L\omega=\omega\circ\mathcal B_L^{-1}$.
    In \cite[lines~825--836]{Nachtergaele1996Spectral}, this pullback is the state transformation associated with the regrouping of
    sites.
    \uses{def:qca_functional_star_equivalence,def:qca_quasi_local_blocking}
\end{definition}

\begin{lemma}[Site grouping preserves states and purity]
    \label{thm:qca_blocking_state_purity}
    \lean{SpinChain.norm_quasiLocalBlockingFunctional,
        SpinChain.quasiLocalBlockingFunctional_mem_stateSpace_iff,
        SpinChain.isPureQuasiLocalState_quasiLocalBlockingFunctional_iff}
    \leanok
    The functional $T_L\omega$ has the same norm as $\omega$, is a state
    exactly when $\omega$ is a state, and is pure exactly when $\omega$ is
    pure. This applies to all state functionals, without an invariance
    condition under translations.
\end{lemma}
\begin{proof}\leanok
    \uses{def:qca_blocking_functional,thm:qca_state_purity_star_equivalence}
    The star-algebra equivalence theorem applied to $\mathcal B_L$ gives all stated properties.
\end{proof}
